\subsection{Station--channel assignment and drainage-area fidelity audit}
\label{sec:methods-drainage-audit}

Objective~2 evaluation includes an intermediate fidelity check that sits between structural preprocessing QA (\S\ref{sec:methods-nhd}) and hydrologic performance: does the \emph{executable} SWAT\textsuperscript{+} model preserve cumulative drainage-area structure relative to the original NHDPlus HR network at gaged channels?
Answering this requires first assigning each USGS streamgage to a SWAT\textsuperscript{+} channel without circularity, then comparing the assigned channel's drainage area to the unmodified hydrography.

\subsubsection*{NHD-first, SWAT-second station assignment}
Gage-to-channel assignment uses a two-stage rule that deliberately excludes SWAT\textsuperscript{+} area from the channel choice, so that the subsequent drainage-area comparison is not tautological.
\emph{NHD-first}: a reference reach is selected on the original NHDPlus HR network from gage coordinates, USGS NWIS site drainage area (where available), GNIS/name context, stream order, LevelPath, \texttt{FType}, and lake/canal flags---never from SWAT\textsuperscript{+} \texttt{AreaC} or \texttt{chandeg.con} area.
\emph{SWAT-second}: that reference reach is mapped to an executable \texttt{chandeg.con} channel by an NHD$\rightarrow$GIS crosswalk, or by a controlled lake-outlet or downstream replacement where the reference reach has no direct executable channel.
Each gage is then labeled with an assignment class (clean tributary, clean mainstem, known mainstem NHD offset, or lake/canal review) and a calibration-eligibility flag.
Drainage-area comparison runs only after assignment.

\subsubsection*{Drainage-area definitions}
The audit distinguishes four drainage-area definitions at each gage: (i)~the \textbf{USGS NWIS site} drainage area from the gage metadata catalog; (ii)~the \textbf{NHD reference} cumulative drainage \texttt{NHDPlusFlowlineVAA.TotDASqKm} on the assigned reach (within 500~m of the gage); (iii)~the \textbf{SWAT\textsuperscript{+}} executable channel area in \texttt{chandeg.con} (hectares $\div$ 100, converted to km\textsuperscript{2}); and (iv)~the \textbf{WBD HU12} upstream-polygon sum, retained as administrative context only.
The NHD reference is read from the domain-clipped USGS HU4 geodatabase---the unmodified source vectors, \emph{not} the cleaned national pickles in which orphan catchments have already been dissolved---so that any offset cannot be an artifact of comparing to a post-processed network.

\subsubsection*{Two alignment comparisons}
Two questions are evaluated separately because they probe different parts of the pipeline.
The \textbf{NHD--USGS} comparison aligns hydrography catalog area (ii) with gage metadata (i), summarized by agreement within 10\% and within 25\% of the NWIS site area; it tests whether the hydrography reach assigned to a gage carries a plausible catalog drainage area.
The \textbf{SWAT\textsuperscript{+}--NHD} comparison aligns the executable model (iii) with the hydrography reference (ii), summarized by the SWAT\textsuperscript{+}/NHD ratio against a loose band ($0.5$--$2.0\times$) and a tight band ($0.8$--$1.25\times$); it tests whether project assembly preserves cumulative drainage structure.
Large SWAT\textsuperscript{+}--NWIS gaps are treated as gage--reach assignment diagnostics, not as evidence of NHD$\rightarrow$SWAT\textsuperscript{+} conversion error.
The SWAT\textsuperscript{+}--NHD audit is applied across the eight catalog Model IDs of the evaluation matrix (Table~\ref{tab:model-roster}); the whole-HUC-8 Peace River domain (\texttt{03100101}, 76 gages) supports per-class analysis of both comparisons.

\subsubsection*{Pipeline-stage decomposition}
Where a systematic SWAT\textsuperscript{+}/NHD offset appears, its origin is localized by tracing cumulative drainage area through successive pipeline stages on the Peace mainstem (Table~\ref{tab:drainage-area-pipeline-trace}): original VAA \texttt{TotDASqKm} (A) and upstream catchment-area sums on the original (B) and cleaned (C) networks, the pre-QSWAT\textsuperscript{+} reach attribute carried in \texttt{streams.pkl} (G), the SWATGenX post-processed watershed-polygon area sum (D), the QSWAT\textsuperscript{+}/TauDEM channel \texttt{AreaC} at assignment (E), and the executable \texttt{chandeg.con} area (F).
A live project channel-area stage (H, \texttt{rivs1.AreaC} and the project SQLite \texttt{gis\_channels} table) confirms whether the offset is already established in the QSWAT\textsuperscript{+} project channel area before the SWAT\textsuperscript{+} text export.
This decomposition attributes any offset to a specific stage rather than to preprocessing in aggregate.
